\documentclass{book}
\usepackage{amsmath, amssymb, amsthm}
\setlength{\parindent}{0pt}

\title{Discrete Structures 1}
\author{Gianna Cantor}
\date{October 30, 2025}

\begin{document}
\maketitle

\chapter{Other Proofs and Methods}

\section{Problem I}
Consider the following proposition about integers in the context of algorithm design:
\begin{center}
    For any integer $n$, $n$ is even if and only if $n^2$ is even.
\end{center}

\noindent \textbf{Proof:} \\
Case 1: $p \rightarrow q$ (Direct Proof)
Assume that $n$ is even. \\
Then $n = 2k$ for some integer $k$. \\
Now, \begin{multline} \label{eq1}
    \begin{split}
        n^2 &= (2k)^2 \\
        n^2 &= 4k^2 \\
        n^2 &= 2(2k)^2
    \end{split}
    \end{multline}
Since $n^2 = 2(integer)$, $n^2$ is even. \\

Case 2: $q \rightarrow p$ (Contraposition)
Assume that $n$ is not even, so $n$ is odd. \\
Then $n = 2k+1$ for some integer $k$. \\
Now, \begin{multline} \label{eq2}
    \begin{split}
        n^2 &= (2k + 1)^2 \\
        n^2 &= (4k^2 + 4k + 1)^2 \\
        n^2 &= 2(2k^2 + 2k + 1)^2
    \end{split}
    \end{multline}
Since $n^2 = 2(integer) + 1$, $n^2$ is odd. \\
However, if $n^2$ is even \eqref{eq1}, then $n$ cannot be odd.
Therefore, $n$ must be even. \\

\vspace{1em}
\fbox{\parbox{\linewidth}{
    $\therefore$ By the biconditional theorem $n$ is even \iff $n^2$ is even, the proposition is true. \textbf{Q.E.D.}
}}

%%%%%%%%%%%%%%%%%%%%%%%%%%%%%%%%%%%%%%
\section{Problem II}
The function $n$ is defined as:
\begin{equation*}
    g(n) = \begin{cases}
        n^2 + 1 &\text{if $n$ is even} \\
        n^2     &\text{if $n$ is odd}
    \end{cases}
\end{equation*}
Claim: For every integer $n$, $g(n)$ is odd. \\

\textbf{Proof:} \\
Case 1: \\
Assume that $n$ is even. \\
Then $n = 2k$ for some integer $k$. \\
Now, \begin{multline} \label{eq3}
    \begin{split}
        n^2 &= (2k)^2 \\
        n^2 &= 4k^2 \\
        n^2 &= 2(2k^2)
    \end{split}
\end{multline}
Since, $n^2 = 2(integer)$, $n^2$ is odd. \\
Therefore, if $n$ is even, then $g(n)$ is odd. \\

Case 2: \\
Assume that $n$ is odd. \\
Then $n = 2k + 1$ for some integer $k$. \\
Now, \begin{multline} \label{eq4}
    \begin{split}
        n^2 &= (2k + 1)^2 \\
        n^2 &= 4k^2 + 4k + 1 \\
        n^2 &= 2(2k^2 + 2k) + 1
    \end{split}
\end{multline}
Since, $n^2 = 2(integer) + 1$, $n^2$ is odd. \\
Therefore, if $n$ is odd, then $g(n)$ is odd. \\

\vspace{1em}
\fbox{\parbox{\linewidth}{
    $\therefore$ For every integer $n$, $g(n)$ is odd if $n$ is even or $n$ is odd. \textbf{Q.E.D.}
}}

%%%%%%%%%%%%%%%%%%%%%%%%%%%%%%%%%%%
\section{Problem III}
Consider the function:
\begin{equation*}
    f(x, y) = (x + y) % 2, where $x$ and $y$ are bits (0 or 1).
\end{equation*}
Claim: $f(x, y) = 1$, if and only if $x \not= y$. \\

\textbf{Proof:} \\
Case 1: \\
Assume $f(x, y) = 1$. \\
Then $f(x, y) = (x + y) \% 2$ \\
Now, \begin{multline} \label{eq5}
    \begin{split}
        (x + y) \% 2 &= (0 + 1) \% 2 &= (1 + 0) \% 2 \\
        (x + y) \% 2 &= (1) \% 2     &= (1) \% 2 \\
        (x + y) \% 2 &= 1            &= 1
    \end{split}
\end{multline}
Therefore, if $f(x, y) = 1$, then $x \not= y$. \\

Case 2: \\
Assume $x \not= y$.
So $x = 0$, $y = 1$ or $x = 1$, $y = 0$.
Now, \begin{multline} \label{eq6}
    \begin{split}
        (x + y) \% 2 &= (0 + 1) \% 2 &= (1 + 0) \% 2 \\
        (x + y) \% 2 &= (1) \% 2     &= (1) \% 2 \\
        (x + y) \% 2 &= 1            &= 1
    \end{split}
\end{multline}
Therefore, if $x \not= y$, then $f(x, y) =1$. \\

\vspace{1em}
\fbox{\parbox{\linewidth}{
    $\therefore$ By the biconditional theorem, the claim is true. $f(x, y) = 1 \iff x \not= y$. \textbf{Q.E.D.}
}}

\end{document}